% !TEX program = xelatex
% 手工维护的 LaTeX（与 keldysh.md 平行；不再经 gen_tex.py/pandoc 生成）。
% 结构：编号恒等式/引理环境 + 短步骤推导 + 表格化约定；更新内容请直接改本文件。
\documentclass[11pt,a4paper]{ctexart}

\usepackage[margin=2.2cm]{geometry}
\usepackage{amsmath,amssymb,mathtools}
\usepackage{amsthm}
\usepackage{booktabs,array}
\usepackage{listings}
\usepackage{xcolor}
\usepackage[colorlinks=true,linkcolor=blue!60!black,urlcolor=blue!60!black]{hyperref}

\setmainfont{DejaVu Serif}
\setmonofont{DejaVu Sans Mono}

\newcommand{\Pf}{\operatorname{Pf}}
\newcommand{\Tr}{\operatorname{Tr}}
\newcommand{\antisym}{\operatorname{antisym}}
\newcommand{\dd}{\mathrm{d}}

\theoremstyle{plain}
\newtheorem{identity}{恒等式}[section]
\newtheorem{lemma}[identity]{引理}
\theoremstyle{definition}
\newtheorem{rulethm}[identity]{规则}
\newtheorem{convention}[identity]{约定}
\newtheorem{checkcert}[identity]{证书}
\theoremstyle{remark}
\newtheorem{remark}[identity]{注}

\lstdefinestyle{julia}{
  basicstyle=\ttfamily\small, keywordstyle=\bfseries,
  commentstyle=\itshape\color{gray}, frame=single, framerule=0.3pt,
  breaklines=true, columns=fullflexible, keepspaces=true,
}

\title{\textbf{严格 IF 实时核理论：Keldysh 轮廓上的逐步推导}\\[2pt]
\large ——bath 积掉、IF Grassmann 变量积掉与 1/2/4 阶 Trotter 组装}
\author{benchmark/exactif 工作笔记（手工维护版）}
\date{2026 年 11 月}

\begin{document}
\maketitle
\tableofcontents
\clearpage

% ============================================================
\section{总览：模型、目标与路线}

\subsection{配套实现与验证入口}

配套实现：\texttt{src/NativeExact.jl}（compact\_delta / sysdynamics\_exact! /
gtlt\_compact）。验证脚本（\texttt{zidentity}、\texttt{verify\_compact}、
\texttt{verify\_gtlt\_align}、\texttt{verify\_align\_u15}、
\texttt{scaling\_compact\_models}）原在 \texttt{validation/}，该目录已清空，
脚本见 git 历史（commit c25ea43 之前）；\texttt{scaling\_compact\_models.jl}
现驻 \texttt{commons/}。

\subsection{本文约定}

\begin{itemize}
\item 每个新变量在首次出现处定义；推导不跳跃——每个等号给理由。
\item 先在最简单情形（单费米模、单 bath 模、单接口）把每条规则显式算出，
      再陈述多模推广及理由。
\item 严格区分两类陈述：\textbf{精确恒等式}（无近似，机器精度成立）与
      \textbf{近似/约定}（Trotter 分割、截断、绝对符号簿记）。
\item 绝对符号（Berezin 积分取向、整体正负）是纯约定，由 Z 恒等式证书与
      Fock 链证书锁定；推导只涉及相对结构与相对符号，并标注"由哪个证书锁定"。
\end{itemize}

\subsection{模型}

杂质模型（单带推导，多带见 \S\ref{sec:multiband}）：
\begin{equation}
H = H_{\mathrm{imp}} + H_q .
\label{eq:model}
\end{equation}

\begin{description}
\item[H\_imp（杂质哈密顿量，非二次型）]
  单带：$H_{\mathrm{imp}} = \epsilon_d\, n$，$n = d^\dagger d$；
  双带 Anderson：$H_{\mathrm{imp}} = \epsilon_d(n_1+n_2) + U\,n_1 n_2$。
  关键性质：$U$ 项是四费米算符，$H_{\mathrm{imp}}$ 一般不是二次型——
  这正是杂质部分必须走精确 Fock 门、不能并入高斯积分的根本原因（\S\ref{sec:schur}）。
\item[$H_q = H_{\mathrm{bath}} + H_{\mathrm{hyb}}$（二次型）]
  $H_{\mathrm{bath}} = \sum_\alpha \epsilon_\alpha c_\alpha^\dagger c_\alpha$；
  $H_{\mathrm{hyb}} = \sum_\alpha (d^\dagger V_\alpha c_\alpha + c_\alpha^\dagger V_\alpha d)$。
  单粒子形式 $H_q = \sum_{ij} h_{ij}\gamma_i^\dagger\gamma_j$
  （$\gamma_0 = d$，$\gamma_\alpha = c_\alpha$），
  \begin{equation}
  h = \begin{pmatrix} 0 & V^\dagger \\ V & H_{\mathrm{bath}} \end{pmatrix},
  \label{eq:hmat}
  \end{equation}
  杂质对角元为 0：杂质自身演化（$\epsilon_d$ 与 $U$）全部归入 $H_{\mathrm{imp}}$。
  该分配使 $h$ 的 bath 侧只含 bath 自由度——是 \S\ref{sec:schur} 高斯积分成立的条件。
\end{description}

\subsection{初态与目标观测量}

初态可分离：bath 热态
$\rho_b = e^{-\beta H_{\mathrm{bath}}}/Z_b$；杂质初态
$\rho_0 = e^{-\beta H_{\mathrm{imp}}}/\operatorname{tr}(e^{-\beta H_{\mathrm{imp}}})$
（任意初态同样支持）。

目标：Keldysh 轮廓 $0\to T$（前向 $+$ 支）$\to 0$（后向 $-$ 支）上的实时观测量：
\begin{itemize}
\item 末态占据数 $n(T)$：$\rho(T) = \operatorname{Tr}_{\mathrm{bath}}[U\rho_{\mathrm{tot}}(0)U^\dagger]$ 的对角元；
\item 双时 Green 函数（裸关联形式；物理 gt/lt 差一个 $\mp i$ 因子）：
      $\mathrm{greater}(k) = \langle c(T)c^\dagger(\tau_k)\rangle$，
      $\mathrm{lesser}(k) = \langle c^\dagger(\tau_k)c(T)\rangle$。
\end{itemize}
全部观测量都是\textbf{比值}：分子 = 带插入收缩，分母 $Z$ = 无插入收缩；
归一化常数（\S\ref{sec:schur} 的 $\nu\Pf$）在比值中消去。

\subsection{总路线}

\begin{table}[htbp]
\centering
\begin{tabular}{cllll}
\toprule
步骤 & 积掉对象 & 方法 & 输出 & 实现位置 \\
\midrule
(1) & bath 相干场 & 逐 cell 精确高斯积分 $+$ Schur & 接口变量上的 Gaussian $G$ & \texttt{compact\_delta} \\
(2) & IF Grassmann 变量 & 带源高斯积分（Pfaffian 恒等式） & 杂质轨迹腿上的消元核 $D$ & \texttt{delta\_eff} 恒等式 \\
(3) & ——（组装） & 调度门链 $+$ 消元核 $+$ 测度/闭合 & 收缩 $\to$ 读出 & \texttt{sysdynamics\_exact!} 等 \\
\bottomrule
\end{tabular}
\caption{三步路线。(1)(2) 是精确恒等式（数值验证：(2) 逐系数 $\le 2.2\times10^{-16}$；
(1) 的 $G$ 与 IFGTEMPO 逐位一致）。(3) 携带全部 Trotter 误差：
order 1/2 全局 $O(\delta t^2)$、order 4 全局 $O(\delta t^4)$（\S\ref{sec:assembly}）。}
\end{table}

记号约定：每条 Grassmann 腿一个生成元；$\bar x$ 型腿记 creation 类、
$x$ 型腿记 annihilation 类；接口（IF）变量对记 $a = (\bar a, a)$，
杂质轨迹变量对记 $x = (\bar x, x)$；多带时带 flavor 下标。

% ============================================================
\section{预备：费米相干态与积分规则}

本节全部在\textbf{单费米模}上手推（Fock 空间只有 $|0\rangle$、
$|1\rangle = d^\dagger|0\rangle$）；多模推广见 \S\ref{sec:multimode}。

\subsection{相干态定义与基本恒等式}

\begin{convention}[相干态（正号约定）]
\begin{equation}
|x\rangle = e^{x d^\dagger}|0\rangle = |0\rangle + x|1\rangle,
\qquad
\langle\bar x| = \langle0| e^{-\bar x d} = \langle0| - \bar x\langle1|.
\label{eq:coh}
\end{equation}
$x, \bar x$ 是外部 Grassmann 生成元：相互反对易、与物理算符 $d, d^\dagger$
反对易、平方为零。指数展开只有两项：$e^{xd^\dagger} = 1 + xd^\dagger$。
\end{convention}

\begin{identity}[ket 本征关系]\label{id:ket}
$d|x\rangle = x\,|x\rangle$.
\end{identity}
\begin{proof}
$d|x\rangle = d|0\rangle + x\,d|1\rangle$；$d|0\rangle = 0$，
$d|1\rangle = d\,d^\dagger|0\rangle = (1 - d^\dagger d)|0\rangle = |0\rangle$，
故 $d|x\rangle = x|0\rangle = x\,|x\rangle$（末步用 $x^2 = 0$）。
\end{proof}

\begin{identity}[bra 本征关系]\label{id:bra}
$\langle\bar x| d^\dagger |0\rangle = -\bar x$，
即 $\langle\bar x|d^\dagger = -\bar x\,\langle\bar x|$ 在 $|0\rangle$ 分量上成立。
\end{identity}
\begin{proof}
$\langle\bar x|d^\dagger|0\rangle = \langle0|d^\dagger|0\rangle - \bar x\langle1|d^\dagger|0\rangle
= 0 - \bar x\cdot 1 = -\bar x$.
\end{proof}

（(2) 的负号与 (1) 的正号是同一约定的两面；二者在 \S\ref{sec:trace} 的迹规则中相互配合。）

\begin{identity}[重叠]\label{id:overlap}
$\langle\bar x|x\rangle = 1 - \bar x\,x \;(= e^{-\bar x x})$.
\end{identity}
\begin{proof}
展开 $(\langle0| - \bar x\langle1|)(|0\rangle + x|1\rangle)$：仅存活项为
$-\bar x\langle1|x\rangle = -\bar x\langle1|1\rangle = -\bar x x$。
\end{proof}

\subsection{分辨率与逐对积分规则}

\begin{rulethm}[逐对积分规则]\label{rule:int}
以 $\int d\bar x\, dx$ 次序：
\begin{center}
\begin{tabular}{lc}
\toprule
单项式 & 积分值 \\
\midrule
$1$          & $0$ \\
$\bar x$     & $0$ \\
$x$          & $0$ \\
$\bar x\,x$  & $\mathbf{-1}$ \\
\bottomrule
\end{tabular}
\end{center}
\end{rulethm}

\begin{rulethm}[分辨率]\label{rule:res}
\begin{equation}
\int d\bar x\, dx\; e^{-\bar x x}\, |x\rangle\langle\bar x| = \mathbb{1}.
\label{eq:res}
\end{equation}
\end{rulethm}

\begin{proof}[验证（规则 \ref{rule:int} 的一致性）]
展开 $e^{-\bar x x} = 1 - \bar x x$，作用于两组基矢，保留积分值非零的项
（恰含一个 $\bar x$ 与一个 $x$）：
\begin{itemize}
\item $|0\rangle \to$ 系数 $-\bar x x$ 积分 $= +1$，还原 $|0\rangle$；
      $|0\rangle\langle1|$ 项：$-x\bar x = +\bar x x$，积分 $= -1$，
      贡献 $+|0\rangle\langle1|$；
\item $|1\rangle \to$ $-x\bar x\,x$ 含两个 $x$ 为零；
\item $|1\rangle\langle1| \to$ $-x\bar x(1 - \bar x x) = +\bar x x$，
      积分 $= -1$，贡献 $+|1\rangle\langle1|$。
\end{itemize}
两组基矢均还原。\textbf{$\int \bar x\,x = -1$ 的负号正是使分辨率为 $+1$ 的约定}；
改用 $dx\,d\bar x$ 次序则全部符号翻转、公式组同样自洽——绝对符号是约定，
相对结构不是。
\end{proof}

\subsection{迹规则与反周期边界}

\begin{rulethm}[迹规则]\label{rule:trace}
\begin{equation}
\Tr[O] = \int d\bar x\, dx\; e^{-\bar x x}\,\langle-\bar x|\,O\,|x\rangle ,
\label{eq:trace}
\end{equation}
bra 端取 $\langle-\bar x|$（\textbf{反周期边界}）：把 bra 定义中 $\bar x \to -\bar x$，
展开的正号给出 $\langle-\bar x| = \langle0| + \bar x\langle1|$。
\end{rulethm}

\begin{checkcert}[{单模验证 $\Tr[e^{\alpha n}] = 1 + e^{\alpha}$}]\label{chk:trace}
$e^{\alpha n}|x\rangle = |0\rangle + e^{\alpha}x|1\rangle$；
$\langle-\bar x|0\rangle = 1$，$\langle-\bar x|1\rangle = +\bar x$。
故 $\langle-\bar x|e^{\alpha n}|x\rangle = 1 - e^{\alpha}\bar x x$；乘测度
$(1 - \bar x x)$：
\begin{equation*}
(1 - \bar x x)(1 - e^{\alpha}\bar x x) = 1 - (1 + e^{\alpha})\bar x x ,
\end{equation*}
积分 $= 1 + e^{\alpha}$，与 Fock 直算一致；$\alpha = 0$ 给
$\Tr[\mathbb{1}] = 2$。
\end{checkcert}

\begin{remark}[为什么边界必须反周期]
迹的对角求和 $\sum_n\langle n|O|n\rangle$ 在相干态表示中把 bra 端 $\bar x$
与 ket 端 $x$ 配对；配对规则要求两变量以 $(\bar x, x)$ 顺序相邻，而 bra 定义
自然给出 $\langle+\bar x|$——换成 $\langle-\bar x|$ 才落到规则 \ref{rule:int}。
这就是"反周期"的全部内容；多模时逐模反周期。
\end{remark}

\subsection{多模推广}\label{sec:multimode}

多模（1 杂质 $+$ $m$ bath）每模一对 $(\bar x, x)$，全体外部生成元相互反对易；
分辨率与迹规则是逐模乘积。分次乘积注意：多模算符串重排的费米子符号 $=$
置换符号；本文涉及多模置换绝对符号处均由 \S\ref{sec:ladder} 的证书锁定，
推导只使用"每生成元在单项式中至多出现一次"的幂零性。

\subsection{高斯矩阵元引理（cell 传播子）}

\begin{lemma}[cell 传播子的相干矩阵元]\label{lem:cell}
对单粒子矩阵 $h$ 与 $U = e^{-i\tau h}$（分块：杂质 $=$ 第 1 行/列，bath $=$ 其余），
二次型算子 $B(\tau) = e^{-i\tau h}$ 的相干态矩阵元是外部变量的高斯型指数，
系数恰为 $U$ 的块：
\begin{equation}
\langle\bar x', \bar a'\,|\,B(\tau)\,|\,x, a\rangle \;\propto\;
\exp\!\big[\,\bar x' U_{11} x + \bar x' U_{12} a + \bar a' U_{21} x + \bar a' U_{22} a\,\big],
\label{eq:celllemma}
\end{equation}
$\propto$ 吸收的常数前因子最终进入 $\Pf(G)$（数值 Pfaffian 计算）。
\end{lemma}

\begin{proof}[单模验证（$V = 0$，$h$ 对角）]
$U = \operatorname{diag}(1, e^{-i\tau\epsilon})$，引理给
$\propto\exp[\bar x' x + e^{-i\tau\epsilon}\bar a' a]$。直接展开：
$e^{-i\tau\epsilon c^\dagger c} = 1 + e^{-i\tau\epsilon}c^\dagger c$（幂零截断），
矩阵元 $= e^{-\bar x'x - \bar a'a} + e^{-i\tau\epsilon}(\bar x' x)(\bar a' a)\cdot(\text{费米符号})$。
两条展开在"每个生成元至多一次"的单项式上一致；常数因子
$e^{-\bar x'x - \bar a'a}$ 正是 \S\ref{sec:res} 型边界分辨率——在链合成中
成为相邻 pair 的测度（\S\ref{sec:qkernel}）。含 $V$ 的单模情形由同一幂零展开
逐系数验证；整体正确性由 \S\ref{sec:ladder} 证书锁定。
\end{proof}

\begin{rulethm}[带源高斯积分]\label{rule:gauss}
\begin{equation}
\int \mathcal{D}a\; e^{\frac12 a^{\mathsf T} G a + j^{\mathsf T} a}
= \Pf(G)\; e^{-\frac12 j^{\mathsf T} G^{-1} j} ,
\label{eq:gauss}
\end{equation}
$a$ = Grassmann 生成元列向量、$G$ 反对称（$G^{\mathsf T} = -G$）、$j$ = 外部源。
$n = 2$ 展开验证见 \S\ref{sec:n2}。
\end{rulethm}

% ============================================================
\section{第一步：积掉 bath（compact\_delta）}

本节对照实现逐步走"积掉 bath"：输入模型参数与调度表
（\texttt{sched.cells} = $N_c$ 个 cell 时长、\texttt{sched.links} = $N_c{+}1$ 个门时长），
输出接口变量上的 Gaussian 核 $G$ 与归一化。实现 $=$ \texttt{NativeExact.compact\_delta}。

\subsection{代码总览：一个大反对称矩阵 \texorpdfstring{$A$}{A}}\label{sec:schur}

全部工作 = 组装 $L\times L$ 反对称矩阵 $A$，再一次 Schur 消元：
\begin{equation}
L = \underbrace{4N_c}_{\text{接口腿}} + \underbrace{2m(2N_c+1)}_{\text{bath 腿}} ,
\end{equation}
$m$ = bath 模式数。代码骨架：
\begin{lstlisting}[style=julia]
nif = 4 * nl                          # 接口腿数（4 条/接口）
total = nif + 2m * (2nl + 1)          # 接口 + bath 腿
A = zeros(ComplexF64, total, total)
h = _coupled_h(V, Hbath)
Us = [exp(-1im * cell .* h) for cell in cells]   # 逐 cell 单粒子传播子
# …… 逐项填充 A（见 2.2-2.3）……
E = A[1:nif, 1:nif]; C = A[1:nif, nif+1:end]; B = A[nif+1:end, nif+1:end]
effective = E + C * (B \ transpose(C))           # Schur：积掉 bath
G = 0.5 .* (effective .- transpose(effective))   # 反对称化
\end{lstlisting}

\subsection{变量布局}

\textbf{接口腿}（前 $4N_c$ 个槽位）：接口 $c$ 的槽位
$4(c{-}1){+}1 \,\ldots\, 4(c{-}1){+}4$ 依次为
$\bar x^{+}(c), x^{+}(c), \bar x^{-}(c), x^{-}(c)$
（代码 \texttt{\_ifslot(c, br, stat)}）。这些腿\textbf{同时}是
(1) 杂质轨迹变量（\S\ref{sec:assembly} 的门链挂在其上）与
(2) Gaussian 核 $G$ 的变量——两扇区记法中它们是两组变量由交叉测度绑定，
紧凑记法直接合一（等价性 \S\ref{sec:fusion}）。

\textbf{bath 腿}：组织成 pair（每个 pair = 一组 bath 相干场
$(\bar\gamma, \gamma)$ 对），共 $2N_c + 1$ 对：

\begin{table}[htbp]
\centering
\begin{tabular}{lll}
\toprule
pair & 个数 & 角色 \\
\midrule
\texttt{forward\_pairs[$c$]}, $c = 1..N_c$ & $N_c$ & forward 链边界：$B(\tau_{c+1})$ 的输出 bars / $B(\tau_c)$ 的输入 etas \\
\texttt{backward\_pairs[$c$]} & $N_c$ & backward 链同构 \\
\texttt{trace\_pair} & $1$ & 链末端的共享 trace pair（\S\ref{sec:qkernel}） \\
\bottomrule
\end{tabular}
\end{table}

\subsection{逐 cell 的 Q-kernel 缝合（四个 U 块）}\label{sec:qkernel}

每个 forward cell $c$ = 一次 \texttt{\_q\_add!}：接口腿
$(\bar x^{+}(c), x^{+}(c))$、输出 bath bars $=$ \texttt{forward\_pairs[c+1].bars}、
输入 bath etas $=$ \texttt{forward\_pairs[c].etas}、传播子 $U^{(c)} = e^{-i\tau_c h}$。
\texttt{\_q\_add!} 放置四个反对称块（\texttt{put\_G!(A,i,j,v)} $=$
$A_{ij} \mathrel{+}= v$ 且 $A_{ji} \mathrel{-}= v$）：

\begin{table}[htbp]
\centering
\begin{tabular}{lll}
\toprule
块 & 系数 & 物理含义 \\
\midrule
接口 pair（bar$_{\mathrm{imp}}$, eta$_{\mathrm{imp}}$） & $U_{11}(\tau_c)$ & \textbf{透射幅}：杂质模式穿过 cell $c$ \\
接口 bar $\times$ bath eta & $U_{1,1+j}$ & \textbf{入射幅}：$d^\dagger V_\alpha$（bath 粒子被吸收） \\
bath bar $\times$ 接口 eta & $U_{1+i,1}$ & \textbf{出射幅}：$V_\alpha d$（发射回 bath） \\
bath $\times$ bath & $U_{22}$ 块 & bath \textbf{内部传播}（$H_{\mathrm{bath}}$ 对角与跃迁） \\
\bottomrule
\end{tabular}
\caption{四个块 $=$ 单粒子传播子 $U$ 的四个象限（引理 \ref{lem:cell}）。
透射幅的诚实陈述：$U_{11}(\tau_c) = [e^{-i\tau_c h}]_{11}$ 一般 $\neq 1$
（含 bath 虚交换的重整化；单模 $= \cos\Omega\tau - i(\epsilon/2\Omega)\sin\Omega\tau$，
$\Omega^2 = \epsilon^2/4 + |V|^2$），$\tau \to 0$ 才趋于 1。}
\end{table}

backward 链：同构 pair 结构，核系数取 $\operatorname{adjoint}(U)$，接口腿用
$\bar x^-/x^-$。

\paragraph{热边} 前/后向链在 $t{=}0$ 由 bath 热态闭合（规则 \ref{rule:trace} 多模版）：
$A_{\bar\gamma^f_i, \gamma^b_j} \mathrel{+}= \mathrm{thermal}[i,j]$，
$\mathrm{thermal} = e^{-\beta H_{\mathrm{bath}}}$。

\paragraph{迹 pair 与 $\langle-\bar\psi|$ 翻转} 最后一个 forward cell 的输出
不是新 pair 而是 \texttt{trace\_pair}（forward 与 backward 链在 $T$ 端共用一对
边界变量，闭合成一个迹），且该核的输出行系数\textbf{翻号}——这正是规则
\ref{rule:trace} 的反周期边界：bra 端生成元整体取负 $=$ 输出行翻号。

\paragraph{逐对测度 $-\bar\gamma\gamma$} 每个 bath pair 上的边界分辨率
（规则 \ref{rule:res} 多模版）；绝对符号 $-1$ 由 Z 证书锁定。

\subsection{Schur 消元：积掉 bath}

至此 $A$ 为（接口 $|$ bath）分块的反对称矩阵，
$\eta_{\mathrm{if}}$ = 接口生成元、$b$ = bath 生成元：
\begin{equation}
A = \begin{pmatrix} E & C \\ -C^{\mathsf T} & B \end{pmatrix},
\qquad
\exp(A) = \tfrac12 \eta_{\mathrm{if}}^{\mathsf T} E\, \eta_{\mathrm{if}}
+ \eta_{\mathrm{if}}^{\mathsf T} C\, b + \tfrac12 b^{\mathsf T} B\, b .
\end{equation}
对 $b$ 用规则 \ref{rule:gauss}（源 $j = C^{\mathsf T}\eta_{\mathrm{if}}$，
二次型系数 $B$；推导 $=$ 配方：平移 $b \to b - B^{-1}C^{\mathsf T}\eta_{\mathrm{if}}$）：
\begin{equation}
\int \mathcal{D}b\;(\cdots) = \Pf(B)\;
e^{\frac12 \eta_{\mathrm{if}}^{\mathsf T}(E + C B^{-1} C^{\mathsf T})\eta_{\mathrm{if}}}.
\label{eq:schur}
\end{equation}
这就是"积掉 bath"：全部 bath Grassmann 变量被精确积掉，剩余多项式只含
接口变量，其 Gaussian 核
\begin{equation}
G = \antisym\!\big(E + C B^{-1} C^{\mathsf T}\big)
= \tfrac12\big(\text{effective} - \text{effective}^{\mathsf T}\big),
\qquad \text{effective} = E + C B^{-1} C^{\mathsf T}.
\label{eq:G}
\end{equation}

\begin{remark}[为什么一次 Schur 就积干净]
bath 变量只通过 $B$ 块相互耦合、通过 $C$ 块与接口耦合；配平后 bath 部分
成为独立 Gaussian，其积分把接口块重整化为 $E + CB^{-1}C^{\mathsf T}$。
多 cell 的链结构（每个 bath pair 出现在相邻两个核中）全部体现在 $C/B$ 的
具体元里，Schur 一次性处理。
\end{remark}

\begin{checkcert}[$G$ 逐位一致]
所得 $G$ 与 IFGTEMPO \texttt{influence\_matrix\_G} 逐位一致
（order 2 的 $4\times4$、order 4 含负 cell 的 $12\times12$，$\max|\Delta G| = 0$）
——pair 布局、热边、迹翻转、测度符号的簿记与 IFGTEMPO 完全相同。
\end{checkcert}

\subsection{U(1) 结构、归一化 \texorpdfstring{$\nu$}{nu} 与小结}

\begin{itemize}
\item \textbf{U(1) 同 kind 零}：$H_q$ 数量守恒 $\Rightarrow$ 四个块全是
      (bar, eta) 跨类对 $\Rightarrow$ $A$ 与 Schur 后的 $G$ 都没有同 kind 元
      （全 order 数值验证为 0）。这是 \S\ref{sec:rows} 行扫表示可行性的关键。
\item \textbf{归一化}：Schur 抛掉的 $\Pf(B)$、逐对测度常数与 bath 热态归一化合成
      $\nu = Z_N/Z$（\texttt{bath\_norm}）。闭合 $Z = (\text{收缩})\cdot\nu\cdot\Pf(G)$；
      比值型观测量中自动消去。
\item \textbf{小结}：$A$ 装配 $+$ 一次 Schur 后，剩余多项式只含接口腿：
      $P(x) = (\text{常数})\,e^{\frac12 x^{\mathsf T}Gx}\,K[x]\cdot(\text{测度/闭合已并入})$。
\end{itemize}

% ============================================================
\section{第二步：积掉 IF Grassmann 变量（消元核 \texorpdfstring{$D$}{D}）}

\subsection{两种记法}

\S\ref{sec:schur} 得到的 $G$ 定义在接口变量上。两扇区记法中接口变量是独立的
IF Grassmann 变量 $a$，与杂质轨迹变量 $x$ 由交叉测度绑定：
\begin{equation}
P(a,x) = \nu\Pf(\mathcal{B})\;
e^{\frac12 a^{\mathsf T} G a}
\prod_c e^{-\bar a_c\, x^{\eta}_c - \bar x^{\bar{}}_c\, a_c}\, K[x]\, M_0 ,
\label{eq:twosector}
\end{equation}
$c$ 遍历接口；$x^{\eta}_c$ = 接口 $c$ 的 annihilation 腿（前向支）、
$\bar x_c$ = creation 腿（后向支）；$M_0$ = $t{=}0$ 分辨率测度。交叉测度是
两扇区构造的接口恒等分辨率（逐接口独立因子——它们的非独立性正是
\S\ref{sec:schur} 的 Schur 要一次做完的原因）。

\textbf{本步}：用规则 \ref{rule:gauss} 把 $a$ 精确积掉（源由交叉测度提供），
积掉后 $a$ 消失、$D$ 落在轨迹腿上。紧凑实现把本步与 \S\ref{sec:schur}
融合为一次大 Schur（等价性 \S\ref{sec:fusion}）。

\subsection{交叉测度作为源（partner 映射）}

\begin{table}[htbp]
\centering
\begin{tabular}{ll}
\toprule
IF 槽位 $s$ & 源 $j_s$（$=$ partner 杂质腿） \\
\midrule
$\bar a_c^{+}$（IF bar，前向） & $-x^{\eta}(c,+)$（杂质 annihilation 腿） \\
$a_c^{+}$（IF eta，前向）       & $-\bar x(c,+)$（杂质 creation 腿） \\
$\bar a_c^{-}$，$a_c^{-}$       & 同构（后向支） \\
\bottomrule
\end{tabular}
\caption{源 $=$ partner 杂质腿（bar $\leftrightarrow$ eta 交叉；
代码 \texttt{partner\_var(t, br, stat) = imp(t, br, flip(stat))}）。}
\end{table}

\subsection{带源高斯积分的 \texorpdfstring{$n=2$}{n=2} 验证}\label{sec:n2}

单接口，$a = (\bar a, a)$，
$G = \begin{pmatrix} 0 & g \\ -g & 0 \end{pmatrix}$（$g$ = 该接口高斯幅），
源 $j^{\mathsf T}a = j_{\bar a}\bar a + j_a a$，$j_{\bar a} = -x^{\eta}$、
$j_a = -\bar x$（$x, \bar x$ 与 $a$ 同属一个 Grassmann 代数——路径积分的机制：
所有腿都是生成元）。逐步展开（全部只到二阶）：
\begin{enumerate}
\item $\tfrac12 a^{\mathsf T}Ga = g\,\bar a\,a$（第二项用 $a\bar a = -\bar a\,a$），
      故 $e^{\frac12 a^{\mathsf T}Ga} = 1 + g\,\bar a\,a$；
\item $e^{j^{\mathsf T}a} = 1 + j_{\bar a}\bar a + j_a a + j_{\bar a}j_a\,\bar a\,a\cdot(+1)$
      （两个源的积 $x\bar x\,\bar a\,a$：$\bar x$ 移过 $\bar a$、$a$ 各反号，净 $(+1)$）；
\item 乘积中含 $\bar a\,a$ 的项之和 $= (g + x\bar x)\,\bar a\,a$
      （4 生成元项为零：$\bar a$ 重复）；
\item 积分（$\int \bar a\,a = +1$）：单接口消元给出
      $g + x\bar x$ —— 杂质腿对上的核 $D = g\,\bar x\,x$ 加常数 $g$；
      绝对符号由证书锁定。
\end{enumerate}
一般 $n$：$e^{j^{\mathsf T}a}$ 展开到二阶，其与 $\tfrac12 a^{\mathsf T}Ga$ 的交叉
把 $-\tfrac12 j^{\mathsf T}G^{-1}j$ 吸收进指数（规则 \ref{rule:gauss}）；杂质腿对的核
由 $G^{-1}$ 的 partner 读出给出（\S\ref{sec:D}）。

\begin{checkcert}[Step B]
本步是精确恒等式：实现层面由暴力 Berezin 积分逐系数对照锁定
（$\le 2.2\times10^{-16}$，orders 1/2/4 全验）。
\end{checkcert}

\subsection{消元核 \texorpdfstring{$D$}{D} 的读出}\label{sec:D}

一般 $N_c$ 接口、$4N_c$ 条杂质腿（每接口 $\bar x^+, x^+, \bar x^-, x^-$）：
\begin{equation}
D[\lambda,\lambda'] = -\tfrac12\big(G^{-1}_{ik} - G^{-1}_{ki}\big),
\qquad i = \operatorname{slot}(\bar\lambda),\ k = \operatorname{slot}(\lambda'),
\label{eq:D}
\end{equation}
$\lambda, \lambda'$ = 两条杂质轨迹腿；$\operatorname{slot}(\bar\lambda)$ =
$\lambda$ 的 partner IF 槽位（bar $\leftrightarrow$ eta 交叉：
$\lambda = \bar x^{+}(c)$ 的 partner 是 $(c, +, \eta)$，
$\lambda' = x^{+}(c')$ 的 partner 是 $(c', +, \mathrm{bar})$）；
$-\tfrac12(\cdot - \cdot)$ = 把 $G^{-1}$ 反对称化后按 bar $\to$ eta 有向读出的
反对称化因子；绝对符号由 Step B 证书锁定（纯约定）。

\subsection{同 kind 零与行分解（Q1 判决）}\label{sec:rows}

\S\ref{sec:schur} 的 U(1) 结构经 Schur 与求逆保持 $\Rightarrow$ $G^{-1}$ 同 kind
元为零 $\Rightarrow$ $D$ 同 kind 元为零（全 order 数值验证 0 个）$\Rightarrow$
\begin{equation}
e^{\sum D} = \prod_{\text{bar } r}\Big(1 + \theta_{\bar r}\sum_c D[r,c]\,\theta_c\Big)
\end{equation}
按行分解——每行一个 partialmpo、键维 2，\texttt{hybriddynamics} 行扫完整承载
（Q1 判决）。消元精确、无近似；代价是 $D$ 非局域（IF 键维按 $4^{N_c}$ 增长，
生产上键维截断）。

\subsection{融合记法与两步记法的等价}\label{sec:fusion}

compact 实现把 \S\ref{sec:schur}（Schur）与本节（Pfaffian 消元）融合为一次
大 Schur：接口槽位直接取为轨迹腿（交叉测度成为 Gaussian 的对角恒等对，
其与轨迹腿的连接由骨架链承担）。两种记法的最终多项式逐位等价——证书链：
(1) $G$ vs IFGTEMPO 逐位；(2) 消元核 vs 暴力积分逐系数；
(3) 管线 vs Fock 离散链机器级；(4) 管线 vs IFGTEMPO 机器级。
两步记法用于推导与验证；融合记法用于实现。

% ============================================================
\section{杂质部分与 1/2/4 阶组装}\label{sec:assembly}

\subsection{杂质核 \texorpdfstring{$K$}{K} 链}

杂质扇区\textbf{不走高斯积分}（$H_{\mathrm{imp}}$ 含 $U$，非二次型）——
走精确 Fock 门：
\begin{equation}
K[x] = \langle \bar x_0^{+}|\rho_0|x_0^{-}\rangle
\prod_{t=1}^{N_c+1}\langle \bar x_t^{+}|A(\ell_t)|x_{t-1}^{+}\rangle
\prod_{t=N_c+1}^{1}\langle \bar x_{t-1}^{-}|A(\ell_t)^{\dagger}|x_t^{-}\rangle,
\qquad A(\ell) = e^{-i\ell H_{\mathrm{imp}}}.
\label{eq:Kchain}
\end{equation}
$A(\ell)$ 对双带含 $U$ 是 $2^{\mathrm{bands}}$ 维 Fock 矩阵的精确指数——
\textbf{门内无近似}；近似只在门与 cell 的分割方式（调度）里。

\subsection{Trotter 调度表}

\begin{table}[htbp]
\centering
\begin{tabular}{clll}
\toprule
order & cells（进 $G$） & links（进 $K$） & 全局阶 \\
\midrule
1 & $\mathrm{fill}(\delta t, N)$ & $[\delta t \times N,\; 0]$ & $O(\delta t^2)$（机理 \S\ref{sec:ord1}） \\
2 & $\mathrm{fill}(\delta t, N)$ & $[\delta t/2,\ \delta t\times(N{-}1),\ \delta t/2]$ & $O(\delta t^2)$（对称） \\
4 & $[a,b,a]\delta t \times N$ & $[a/2,\ c,\ c,\ (a,c,c)\times(N{-}1),\ a/2]$ & $O(\delta t^4)$ \\
\bottomrule
\end{tabular}
\caption{调度表。links 比 cells 多一个（链两端各一门）。
order 1 links 末位 0 的恒等门不可跳（\S\ref{sec:conv}）。}
\end{table}

\subsection{回文与阶条件（order 2 与 4 的推导）}

记 $X = -i\tau_1 A$、$Y = -i\tau_2 B$、$Z = -i\tau_3 A$。对
$e^X e^Y e^Z$，BCH 展开的指数二阶项：
\begin{equation}
\tfrac12[X,Y] + \tfrac12[Y,Z] + \tfrac14\big([X,Z] + [Z,X]\big).
\label{eq:bch2}
\end{equation}

\paragraph{order 2}
链 $X = -i\delta t A/2$、$Y = -i\delta t B$、$Z = -i\delta t A/2$：
$[X, Z] = 0$（同为 $A$ 向），且 $[Y,Z] = [Y,X]$（$\tau$ 相等），于是
$\tfrac12[X,Y] + \tfrac12[Y,X] = 0$——$\delta t^2$ 项恰好相消
（"对称/回文"的代数含义：首尾两个 A-半步相等使两个 commutator 互抵）。
残留从 $\delta t^3$ 起：单步 $O(\delta t^3)$、$N$ 步全局 $O(\delta t^2)$。

\paragraph{order 4}
回文 Suzuki $S_4 = S_2(a\delta t)\,S_2(b\delta t)\,S_2(a\delta t)$，合并链后
逐宏步
\begin{equation}
A(a/2)\,B(a)\,A(c)\,B(b)\,A(c)\,B(a)\,A(a/2),
\qquad c = \frac{a + b}{2},
\end{equation}
权重条件
\begin{equation}
2a + b = 1\ \ (\text{一阶：总时长})，\qquad
2a^3 + b^3 = 0\ \ (\text{三阶相消})，
\end{equation}
解出 $a = 1/(2 - \sqrt[3]{2}) \approx 1.2252$，
$b = 1 - 2a \approx -1.4505 < 0$。

\begin{remark}[回文性的角色与负 cell 的必然性]
组合中 S2 因子序正读倒读一致，所有奇幂 BCH 余项成对相消；非回文序
$S_2(a)^2 S_2(b)$ 的 $\delta t^3$ 项不满足该条件——实测比值序列
$12.7 \to 9.0$（滑向 8，即 $p \to 3$）是其失效的数值表现；修正为回文序后
比值恢复 16（$N{=}16$ 定点 15.97/16.12）。负 cell 的必然性：
$2a^3 + b^3 = 0$ 要求 $b < 0$——三阶相消必须有反向步，接口格点因此
非均匀且时间倒卷。
\end{remark}

\subsection{order 1 全局二阶的机理}\label{sec:ord1}

order 1 的裸分裂 $A(\delta t)B(\delta t)$ 单步只有 $O(\delta t)$——但严格核
管线的实测全局阶是 $p = 2$。机理：分裂的领头误差 $[A,B]\delta t^2/2$ 是
"杂质门与耦合演化的交错缺失"；严格 cell $B(\delta t) = e^{-i\delta t\,h}$ 的
耦合顶点幅对 $V$ 的任意幂精确，它把这项交错重新生成到所有阶（$D$ 的二阶
交叉项正是 $[A,B]\delta t^2$ 的内容），残留从更高阶起，全局 $O(\delta t^2)$。
冻结核没有这个重求和，全局 $O(\delta t)$（实测 $p \approx 1$）。

\subsection{组装约定逐条}\label{sec:conv}

\begin{table}[htbp]
\centering
\begin{tabular}{lll}
\toprule
约定 & 内容 & 理由 \\
\midrule
每 gap 一门 & $\tau = 0$ 恒等门不可跳 & 恒等传播子携带分辨率 $\prod(1 + \bar c c)$，是测度的一部分 \\
接口反测度 & $\mathrm{cc} = -1$ & 幂零代数 $(1 - \bar x x)(1 + \bar x x) = 1$ 精确抵消 \\
归一化 & $\nu\Pf(G)$ 进 scaling & 闭合 $Z$；比值型观测量自动消去 \\
消元符号 & $\mathrm{sgn} = -1$ & Step B 证书锁定 \\
无 siteflip & $M$ 不预翻 & partialmpo 内部处理位点序 \\
门符号 & 前向/后向共轭 & Keldysh 两支方向 \\
\bottomrule
\end{tabular}
\end{table}

\subsection{读出约定与多带推广}\label{sec:multiband}

\begin{itemize}
\item $n(T)$：折返点自由腿对 GTerm 投影。
\item 双占据：双折返腿 GTerm 乘积。
\item gt/lt：$c(t_f)$ @ backward-eta $+$ GTerm 对（有效半步时刻）。
\item 多带：IF 逐带复制；门用多带 Fock 传播子。
\end{itemize}

% ============================================================
\section{闭环验证阶梯（每个数字的出处）}\label{sec:ladder}

\begin{table}[htbp]
\centering
\begin{tabular}{llll}
\toprule
阶梯 & 判据 & 结果 & 出处 \\
\midrule
恒等式级 & $\epsilon_d = 0$: $Z = 1$ & $4.6\times10^{-13}$ & zidentity \\
代数级 & $\int\mathcal{D}a = \nu\Pf\cdot e^{D}$ & $\le 2.2\times10^{-16}$ & Step B \\
离散链级 & vs 同调度 Fock 直算 & $\le 3.8\times10^{-15}$ & verify\_compact \\
对齐级 & vs IFGTEMPO（U=1.5 双带） & $\le 8.6\times10^{-14}$ & verify\_align\_u15 \\
连续极限级 & vs ED，order 1/2/4 & $p = 2/2/4$ & scaling\_compact\_models \\
\bottomrule
\end{tabular}
\caption{验证阶梯（脚本见 git 历史；\texttt{scaling\_compact\_models} 现驻 \texttt{commons/}）。}
\end{table}

% ============================================================
\section{边界与延拓}

\begin{itemize}
\item quapi 路线：仅 orders 1/2。
\item 连续谱 bath：order 4 非均匀 cells 显式拒绝。
\item IFGTEMPO：验证参照。
\end{itemize}

\end{document}
